% Ad Atticum 14.21 (ad-atticum-14-21) --- English translation
% Translated by Claude (Anthropic) from the Latin (Perseus).
% Style guide: STYLE.md.

\ciceroLetterOpener{CICERO ATTICO salutem}

\ciceroSection{1} When I had just a little while ago handed a letter to Cassius's courier for you, on the fifth before the Ides our own courier turned up --- and, like a portent, without a letter from you. I quickly guessed, though, that you had been at Lanuvium. Eros, however, made haste, so that a letter of Dolabella's might be brought through to me --- not on my business (for he had not yet had mine), but a reply, written really very handsomely, to those of which I had sent you the copy.

\ciceroSection{2} As for me, no sooner had I dismissed Cassius's courier than there was Balbus. Good gods, how easily you could see that he feared peace! And you know how guarded the man is. Still, he was full of Antony's plans: that he was going round to the veterans to get them to ratify Caesar's acts, and to swear that they themselves would uphold them; that all of them should keep camp, and that two-man commissioners should inspect them every month. He even complained of the unpopularity falling upon himself; and his whole talk was such that he seemed to love Antony. In a word: not a sincere syllable.

\ciceroSection{3} But, for my part, I have no doubt that the whole thing is moving toward armed camps. For that business was carried through with the spirit of men but the planning of children. For who has not seen this --- that the heir of the kingship has been left in being? What more absurd than this --- \textit{to fear the one, and not to put the other in fear}? Even at this very moment a great deal is \textit{quite a bit barbarous} (\textit{[Greek: huposolokia]}). That Pontius's house at Naples is in the possession of the tyrannicide's mother! I had better read more often the \textit{Cato the Elder} I sent you. For old age makes me more sour. I take everything to heart. But as for me, my life has been lived (\textit{[Greek: bebi\=otai]}); let the young see to it. You will see to my affairs, as you are doing.

\ciceroSection{4} I wrote this --- or rather dictated it --- with the dessert course set out at Vestorius's. Tomorrow I was thinking of dining at Hirtius's --- and, indeed, a five-course meal (\textit{[Greek: pentelopion]}). So I am preparing to bring the man across to the side of the \textit{optimates}. Tall talk (\textit{[Greek: l\=eros polus]}). There is not one of that lot who does not fear peace. So let us look about for our travelling-shoes. Anything sooner than the camps. Give Attica my warmest love, please. I am waiting for Octavian's harangue, and for anything else there is --- but most of all whether Dolabella jingles a bit, or has cancelled his debts under my name.
